\documentclass[10pt,a4paper]{article}
\usepackage[T1]{fontenc}
\usepackage[utf8]{inputenc}
\usepackage[brazil]{babel}
\usepackage{lmodern}
\usepackage{microtype}
\usepackage{geometry}
\usepackage[table]{xcolor}
\usepackage{longtable}
\usepackage{array}
\usepackage{booktabs}
\usepackage{ragged2e}
\usepackage{fancyhdr}
\usepackage{calc}
\usepackage{lastpage}
\newcolumntype{P}[1]{>{\RaggedRight\arraybackslash}p{#1}}
\setlength{\parindent}{0pt}
\setlength{\parskip}{0pt}
\setlength{\LTpre}{0pt}
\setlength{\LTpost}{0pt}
\setlength{\emergencystretch}{3em}
\geometry{a4paper,landscape,margin=17mm}
\renewcommand{\familydefault}{\sfdefault}
\pagestyle{fancy}
\fancyhf{}
\renewcommand{\headrulewidth}{0pt}
\fancyfoot[L]{\fontsize{7}{9}\selectfont PROTÓTIPO VISUAL \textbullet\ NÃO É A TABELA FINAL DO TCC}
\fancyfoot[R]{\fontsize{7}{9}\selectfont \thepage}
\setlength{\headheight}{13pt}
\begin{document}
{\small\bfseries AMOSTRAS DE SÍNTESE BIBLIOGRÁFICA}\par
\vspace{2mm}
{\Large 03 \textemdash\ Eixos temáticos / síntese organizada}\par
\vspace{2mm}
{\footnotesize Conteúdo ilustrativo para comparar diagramação. Confirmar dados e referências no TCC antes de usar.}\par
\vspace{4mm}
\definecolor{forest}{HTML}{1E4943}
\definecolor{softforest}{HTML}{E8F0EE}
\definecolor{lightline}{HTML}{C8D8D4}
\setlength{\tabcolsep}{5.5pt}
\renewcommand{\arraystretch}{1.55}
\footnotesize
\begin{longtable}{@{}P{0.130\textwidth}P{0.150\textwidth}P{0.190\textwidth}P{0.240\textwidth}P{0.220\textwidth}@{}}
\caption{Síntese temática da fundamentação e dos estudos relacionados}\label{tab:amostra_eixos}\\
\arrayrulecolor{forest}\toprule[1.6pt]
\color{forest}\bfseries Autor(es) & \color{forest}\bfseries Foco do estudo & \color{forest}\bfseries Método / contexto & \color{forest}\bfseries Principais contribuições & \color{forest}\bfseries Uso no presente estudo\\
\arrayrulecolor{lightline}\midrule[0.5pt]
\endfirsthead
\multicolumn{5}{@{}l}{\footnotesize\itshape Tabela \thetable\ (continuação)}\\[1mm]
\arrayrulecolor{forest}\toprule[1.6pt]
\color{forest}\bfseries Autor(es) & \color{forest}\bfseries Foco do estudo & \color{forest}\bfseries Método / contexto & \color{forest}\bfseries Principais contribuições & \color{forest}\bfseries Uso no presente estudo\\
\arrayrulecolor{lightline}\midrule[0.5pt]
\endhead
\multicolumn{5}{r@{}}{\footnotesize Continua na página seguinte}\\
\endfoot
\arrayrulecolor{forest}\bottomrule[1pt]
\multicolumn{5}{@{}l}{\footnotesize\itshape Fonte: elaboração própria (amostra demonstrativa).}\\
\endlastfoot
\rowcolor{softforest}\multicolumn{5}{@{}l}{\color{forest}\bfseries\footnotesize FUNDAMENTOS DE FINANÇAS}\\[1pt]
\textbf{Markowitz (1952)} &
Seleção de portfólios &
Modelo média-variância; retorno esperado e covariância. &
Formaliza o equilíbrio entre risco e retorno na alocação de ativos. &
Base conceitual da otimização de carteiras. \\
\arrayrulecolor{lightline}\specialrule{0.4pt}{0pt}{0pt}
\textbf{Sharpe (1964)} &
Risco e retorno de ativos &
CAPM e risco sistemático de mercado. &
Relaciona o retorno esperado à exposição ao risco sistemático. &
Contextualiza a teoria financeira de alocação. \\
\arrayrulecolor{lightline}\specialrule{0.4pt}{0pt}{0pt}
\rowcolor{softforest}\multicolumn{5}{@{}l}{\color{forest}\bfseries\footnotesize ML APLICADO A FINANÇAS}\\[1pt]
\textbf{White (1988)} &
Previsão de retornos &
Rede neural feedforward; retornos diários de ações da IBM. &
Marco inicial da investigação de redes neurais em séries financeiras. &
Contextualização histórica das aplicações de ML. \\
\arrayrulecolor{lightline}\specialrule{0.4pt}{0pt}{0pt}
\textbf{Wang et al. (2020)} &
Pré-seleção de ativos &
LSTM para previsão e otimização média-variância. &
Integra informação preditiva à escolha de ativos e à alocação. &
Exemplo de arquitetura híbrida para carteiras. \\
\arrayrulecolor{lightline}\specialrule{0.4pt}{0pt}{0pt}
\rowcolor{softforest}\multicolumn{5}{@{}l}{\color{forest}\bfseries\footnotesize ESTUDOS RELACIONADOS}\\[1pt]
\textbf{Ma et al. (2021)} &
Previsão e otimização &
Modelos de ML e DL; otimização MV e Ômega. &
Compara previsores usados em estratégias de alocação de carteiras. &
Antecedente direto da combinação ML e MV. \\
\arrayrulecolor{lightline}\specialrule{0.4pt}{0pt}{0pt}
\textbf{Jiang et al. (2024)} &
Decisão de portfólio &
Aprendizado por reforço profundo; seleção de carteiras. &
Explora políticas de alocação aprendidas a partir de dados. &
Contraponto às abordagens previsão + otimização. \\
\arrayrulecolor{lightline}\specialrule{0.4pt}{0pt}{0pt}
\textbf{Chen e Zheng (2025)} &
Previsão de criptoativos &
LSTM e rede multitarefa; mercado de criptomoedas. &
Investiga previsão e definição de carteiras com redes neurais. &
Comparação de alternativas de ML para criptoativos. \\
\arrayrulecolor{lightline}\specialrule{0.4pt}{0pt}{0pt}
\textbf{Peykani et al. (2026)} &
Risco e diversificação &
XGBoost, otimização bayesiana e estratégias de carteira. &
Compara alternativas de otimização com distintas medidas de risco. &
Relaciona estimação por ML e otimização de portfólios. \\
\arrayrulecolor{lightline}\specialrule{0.4pt}{0pt}{0pt}
\end{longtable}
\end{document}
